\section{结论}
本文提出了一种面向中文文献的图神经网络引用预测框架 Papex-GNN，通过预训练语义编码与异构图注意力结构建模的有机结合，在中文文献引用预测任务上取得了显著优于现有基线的效果。消融实验验证了两类信号各自的不可或缺性。

未来的工作将探索：(1) 引入会议、期刊等更丰富的节点类型以进一步丰富异构图结构；(2) 将框架扩展至多语言文献场景；(3) 结合大语言模型进行可解释的引用推荐。我们相信，语义与结构双轮驱动的建模思路，能够为学术知识图谱的构建与应用提供持续的动力。
